% Part: first-order-logic
% Chapter: proof-systems
% Section: seqeunt-calculus

\documentclass[../../../include/open-logic-section]{subfiles}

\begin{document}

\iftag{FOL}
      {\olfileid{fol}{prf}{seq}}
      {\olfileid{pl}{prf}{seq}}

\olsection{حساب دنباله‌ای}

بسیاری از دستگاه‌های !!{derivation} با آرایش‌هایی از !!{sentence}s
کار می‌کنند، اما حساب دنباله‌ای با \emph{دنباله‌ها} سروکار دارد. دنباله
عبارتی است به‌صورت
\[
!A_1, \dots, !A_m \Sequent !B_1, \dots, !B_m,
\]
یعنی زوجی از توالی‌های !!{sentence}s که نماد دنبالهٔ~$\Sequent$ آن‌ها
را از هم جدا می‌کند. هر یک از دو توالی ممکن است تهی باشد.
!!^a{derivation} در حساب دنباله‌ای درختی از دنباله‌هاست که در آن
بالاترین دنباله‌ها صورت خاصی دارند (آن‌ها را «دنباله‌های آغازین» یا
«اصول موضوع» می‌نامند) و هر دنبالهٔ دیگر با یکی از قواعد استنتاج از
دنباله‌های بی‌درنگ بالای خود به دست می‌آید. قواعد استنتاج یا
!!{sentence}s درون دنباله‌ها را دست‌کاری می‌کنند (آن‌ها را در سمت چپ
یا راست می‌افزایند، حذف می‌کنند یا از نو می‌آرایند)، یا !!{formula}
مرکبی را در نتیجهٔ قاعده وارد می‌کنند. برای مثال، قاعدهٔ
$\LeftR{\land}$ استنتاج از $!A, \Gamma \Sequent \Delta$ به
$!A \land !B, \Gamma \Sequent \Delta$ را مجاز می‌کند، و قاعدهٔ
$\RightR{\lif}$ استنتاج از $!A, \Gamma \Sequent \Delta, !B$ به
$\Gamma \Sequent \Delta, !A \lif !B$ را، برای هر $\Gamma$، $\Delta$،
$!A$ و~$!B$، مجاز می‌کند. (به‌ویژه، $\Gamma$ و~$\Delta$ می‌توانند
تهی باشند.)

رابطهٔ $\Proves$ مبتنی بر حساب دنباله‌ای چنین تعریف می‌شود:
$\Gamma \Proves !A$ اگر و تنها اگر توالی‌ای چون $\Gamma_0$ وجود داشته
باشد که هر $!A$ در $\Gamma_0$ عضوی از~$\Gamma$ باشد و !!{derivation}
با دنبالهٔ~$\Gamma_0 \Sequent !A$ در ریشهٔ خود وجود داشته باشد. $!A$
در حساب دنباله‌ای قضیه است اگر دنبالهٔ~$\Sequent !A$ دارای
!!a{derivation} باشد. برای مثال، !!a{derivation} زیر نشان می‌دهد که
$\Proves (!A \land !B) \lif !A$:
\begin{prooftree}
  \Axiom$!A \fCenter !A$
  \RightLabel{\LeftR{\land}}
  \UnaryInf$!A \land !B \fCenter !A$
  \RightLabel{\RightR{\lif}}
  \UnaryInf$\fCenter (!A \land !B) \lif !A$
\end{prooftree}

مجموعهٔ $\Gamma$ در حساب دنباله‌ای ناسازگار است اگر !!a{derivation}
از $\Gamma_0 \Sequent$ وجود داشته باشد (به‌طوری‌که هر
$!A \in \Gamma_0$ عضوی از~$\Gamma$ باشد و سمت راست دنباله تهی باشد).
با استفاده از قاعدهٔ \RightR{\Weakening} می‌توان هر !!{sentence} را
از مجموعه‌ای ناسازگار !!{derive} کرد.

حساب دنباله‌ای را گرهارد گنتسن در دههٔ ۱۹۳۰ ابداع کرد. طراحی
نظام‌مند و متقارن آن، این حساب را به صورت‌گرایی بسیار سودمند برای
پرورش نظریه‌ای دربارهٔ !!{derivation}s بدل می‌کند. یافتن
!!{derivation}s در حساب دنباله‌ای نسبتاً آسان است، اما خواندن این
!!{derivation}s اغلب دشوار است و پیوند آن‌ها با برهان‌ها گاه به‌آسانی
دیده نمی‌شود. بااین‌حال، حساب دنباله‌ای رویکردی بسیار ظریف به
دستگاه‌های !!{derivation} از کار درآمده است و بسیاری از منطق‌ها
دستگاه‌هایی از نوع حساب دنباله‌ای دارند.

\end{document}
